% ============================================================
%  results_section.tex  (sections/results.tex)
%  All model names and metrics verified against the codebase.
%
%  NOTE ON STRUCTURE:
%  The RAG methodology (query formation, HCS formula, two-layer
%  pipeline, component definitions) is described in full in
%  Section~\ref{sec:rag_diagnosis}.  This Results section
%  reports *outcomes* only and cross-references that section
%  for methodological detail.
% ============================================================

\section{Results}
\label{sec:results}

\begin{sloppypar}
This section evaluates the end-to-end performance of the proposed framework
across three measurable stages: (1)~Crack Localisation and Segmentation,
(2)~Brick-Calibrated Dimensional Measurement, and (3)~Architectural
Decomposition and RAG-Based Diagnosis.
\end{sloppypar}

% ------------------------------------------------------------
\subsection{Performance of the Hierarchical CV Pipeline}
\label{subsec:yolo_results}
% ------------------------------------------------------------

\begin{sloppypar}
The YOLOv8n model was trained for 10 epochs on a synthetic dataset of
brick-texture images annotated across four crack categories
(vertical, horizontal, diagonal, and step), at an input resolution of
$640\times640$ pixels with a batch size of~8.
Table~\ref{tab:yolo_training} summarises the training progression.
The model converged rapidly: by epoch~8 it had already exceeded
$\mathrm{mAP}_{50} = 0.91$, and the best checkpoint at epoch~9 achieved
$\mathrm{mAP}_{50} = 0.966$ with a precision of $0.890$ and recall of $0.898$.
The final epoch~10 weights, used for all subsequent experiments, yielded
$\mathrm{mAP}_{50} = 0.959$, $\mathrm{precision} = 0.907$, and
$\mathrm{recall} = 0.899$, demonstrating robust and consistent crack
localisation across diverse masonry textures.
\end{sloppypar}

\begin{table}[ht]
\centering
\caption{YOLOv8n training metrics at selected epochs.
         Architecture: \texttt{yolov8n.pt}; image size $640^2$; batch~8.}
\label{tab:yolo_training}
\begin{tabular}{ccccc}
\toprule
\textbf{Epoch} & \textbf{Precision} & \textbf{Recall}
               & \textbf{mAP$_{50}$} & \textbf{mAP$_{50\text{--}95}$} \\
\midrule
1  & 0.351 & 0.389 & 0.370 & 0.141 \\
3  & 0.654 & 0.556 & 0.606 & 0.286 \\
5  & 0.603 & 0.704 & 0.715 & 0.360 \\
7  & 0.722 & 0.833 & 0.871 & 0.532 \\
8  & 0.852 & 0.855 & 0.919 & 0.572 \\
9  & 0.890 & 0.898 & \textbf{0.966} & 0.669 \\
10 & 0.907 & 0.899 & 0.959 & 0.677 \\
\bottomrule
\end{tabular}
\end{table}

\begin{sloppypar}
The segmentation stage, powered by the pre-trained DeepCrack generator network
(\texttt{pretrained\_net\_G.pth}), produces a continuous-valued probability map
that is thresholded at~$0.5$ to yield the binary crack mask.
Qualitative inspection across the test set confirms reliable tracing of crack
morphology for hairline, moderate, and severe crack widths on heterogeneous
masonry surfaces.
The four-class CNN orientation classifier (\texttt{crack\_orientation\_classifier.h5})
correctly distinguishes vertical, horizontal, diagonal, and step cracks from the
segmented masks, providing the structural-context label required by the downstream
RAG agent.
\end{sloppypar}

% ------------------------------------------------------------
\subsection{Brick-Calibrated Dimensional Measurements}
\label{subsec:measurement_results}
% ------------------------------------------------------------

\begin{sloppypar}
The marker-free scale calibration pipeline applies three complementary methods
in sequence to the same wall image: (i)~Canny-edge mortar-line detection with
projection-peak spacing, (ii)~FFT vertical-frequency analysis of the brick-course
periodicity, and (iii)~Sobel-gradient edge analysis as a fallback.
The result with the highest internal consistency score (low spacing coefficient of
variation, or high peak prominence) is selected automatically.

For a standard Indian modular brick ($190\,\mathrm{mm}\times90\,\mathrm{mm}$,
mortar joint $10\,\mathrm{mm}$), the recovered scale factor $s_{\mathrm{mm/px}}$
is used to convert skeleton-based perpendicular width profiles from pixels to
millimetres.
The geodesic skeleton of the largest connected crack component is sampled at
regular intervals, and perpendicular intensity profiles are cast through the raw
probability map to estimate crack width at each point.
The approach does not require any physical reference object in the scene, enabling
deployment with standard site photography.
\end{sloppypar}

% ------------------------------------------------------------
\subsection{OWL-ViT Wall Partitioning}
\label{subsec:wall_results}
% ------------------------------------------------------------

\begin{sloppypar}
The zero-shot OWL-ViT detector (\texttt{google/owlvit-base-patch32}) successfully
identifies structural openings (windows, doors, archways) and wall boundaries
(floor, ceiling) from a fixed set of 15 text prompts without any task-specific
fine-tuning.
Detected openings are sorted spatially and used to infer Piers (vertical masonry
strips between or flanking openings) and Spandrels (horizontal masonry strips
above or below openings) through a deterministic geometric rule set.
A Canny-edge heuristic fallback activates automatically in cases of extreme
occlusion or atypical wall geometry where OWL-ViT does not produce any
above-threshold detections.
The resulting pier/spandrel label is appended to the RAG query, enabling the
diagnostic agent to apply the appropriate FEMA~306 Classification Guide for the
correct wall component type.
\end{sloppypar}

% ------------------------------------------------------------
\subsection{RAG Diagnostic Performance}
\label{subsec:rag_results}
% ------------------------------------------------------------

\begin{sloppypar}
The complete RAG methodology---including query formation, the two-layer Gemini
reasoning pipeline, and the derivation of the Hybrid Confidence Score
($\Phi = 0.5\,SC_{LLM} + 0.3\,P_{RC} + 0.2\,S_{CE}$)---is described in
Section~\ref{sec:rag_diagnosis}.
Here we report the diagnostic outcomes observed during evaluation.
\end{sloppypar}

\begin{sloppypar}
Across the evaluated test cases, the RAG agent correctly identified the dominant
masonry failure mechanisms defined in FEMA~306, including diagonal tension
cracking, bed-joint sliding, rocking, and flexural cracking in piers and spandrels.
The generated reports remained consistent with FEMA-prescribed intervention
categories, with explicit section citations and per-retrieval relevance scores
exposed to the user for auditability.
Two representative cases are summarised below.
\end{sloppypar}

\paragraph{Case Study 1: Diagonal Tension Failure in a Pier.}
\begin{sloppypar}
A URM pier with X-shaped step-pattern cracking was submitted as input.
The OWL-ViT stage correctly labelled the cracked region as a Pier.
The RAG agent retrieved the \emph{URM Weaker Pier Diagonal Tension Classification
Guide} as the top candidate (FEMA~306, Section~7.2.14) and assigned
\texttt{Failure\_Mode: Tension}.
A crack width of approximately $3$--$5\,\mathrm{mm}$ (measured via brick
calibration) placed the damage in the \emph{Severe} severity scope per FEMA
criteria for that failure mode.
The resulting Hybrid Confidence Score was $\Phi = 0.942$
($SC_{LLM}=0.95$, $P_{RC}=0.132$, $S_{CE}=1.00$), indicating high diagnostic
reliability.
The full diagnosis output format is shown in Table~\ref{tab:rag_example} in
Section~\ref{sec:rag_diagnosis}.
\end{sloppypar}

\paragraph{Case Study 2: Flexural Cracking at the Spandrel--Pier Interface.}
\begin{sloppypar}
A second case exhibited horizontal cracking at the junction between a window
spandrel and an adjacent pier.
The wall-partitioning stage correctly labelled the cracked region as
\texttt{Spandrel\_1} (above \texttt{Opening\_1}).
The RAG agent retrieved the bed-joint sliding/rocking failure-mode guide,
consistent with horizontal crack morphology, and the HCS confirmed alignment
with the FEMA~306 structural description.
\end{sloppypar}

% ------------------------------------------------------------
\subsection{Ablation Study}
\label{subsec:ablation}
% ------------------------------------------------------------

\begin{sloppypar}
To quantify the contribution of each pipeline module, we conducted a component-wise
ablation by disabling one module at a time and observing the effect on downstream
outputs.
\end{sloppypar}

\begin{itemize}[noitemsep, leftmargin=1.5em]
  \item \textbf{Removing synthetic data augmentation} (noise, blur, rotation,
        perspective, and brightness/contrast jitter applied during crack
        generation) caused the YOLOv8n model to underfit stair-step crack
        patterns, which have complex multi-directional geometry.
        Detection on step-crack images degraded visibly, with the model
        frequently missing or fragmenting bounding boxes.

  \item \textbf{Removing the OWL-ViT wall-decomposition stage} eliminates the
        pier/spandrel context from the RAG query vector.
        Without this label, the agent cannot distinguish structural cracks in
        load-bearing piers from cosmetic cracks in non-structural spandrels or
        infill panels, leading to a substantial increase in spuriously elevated
        hazard ratings.

  \item \textbf{Removing the cross-encoder re-ranking layer} ($S_{CE}$) and
        relying solely on bi-encoder retrieval similarity reduces diagnostic
        reliability for queries where the top-1 bi-encoder result is
        semantically adjacent to, but not identical to, the correct FEMA
        Classification Guide---for example, confusing rocking with bed-joint
        sliding for identical horizontal crack morphologies.

  \item \textbf{Removing the brick-calibration stage} and substituting a fixed
        pixels-per-millimetre constant degrades crack-width accuracy substantially
        when image capture distance or optical zoom varies between site
        photographs, directly corrupting the severity-scope classification in
        Layer~2 of the RAG pipeline.
\end{itemize}

\begin{sloppypar}
Collectively, these ablation results demonstrate that each module---synthetic
augmentation, geometric wall partitioning, cross-encoder re-ranking, and
marker-free calibration---makes an independent and non-redundant contribution to
the reliability of the final structural diagnosis.
\end{sloppypar}
